\documentclass[12pt,a4paper]{article}
\usepackage[UTF8]{ctex}
\usepackage{geometry}
\usepackage{fancyhdr}
\usepackage{setspace}

\geometry{left=2.5cm,right=2.5cm,top=2.5cm,bottom=2.5cm}
\setlength{\headheight}{15pt}
\setstretch{1.5}

\pagestyle{fancy}
\fancyhf{}
\fancyhead[C]{说明书摘要}
\fancyfoot[C]{\thepage}

\begin{document}

\begin{center}
    {\Huge\heiti 说明书摘要}
\end{center}

\vspace{1em}

\noindent\textbf{发明名称：}一种面向6G网络切片的超低时延抗量子自适应握手方法、系统、设备及存储介质

\vspace{1em}

\noindent 本发明涉及移动通信安全和后量子密码学技术领域，公开了一种面向6G网络切片的超低时延抗量子自适应握手方法、系统、设备及存储介质。该方法包括：接收携带切片标识和终端能力信息的握手请求；根据切片握手画像、终端能力信息和链路状态信息选择握手计划；执行基础后量子密钥协商以获得基础份额并导出基础会话密钥；在允许早期业务的情况下导出早期业务密钥；生成增强密钥材料并按照分片参数编码为多个增强片段；在增强阶段时限内发送或接收所述增强片段，在达到恢复阈值时恢复增强密钥材料；基于基础份额和增强密钥材料导出最终会话密钥；将切片标识、握手计划标识、算法包摘要和分片参数绑定到握手转录文本中，以实现抗降级和抗跨切片重放保护。本发明能够根据不同网络切片的时延预算、安全等级和终端能力动态拆分与调度后量子密钥材料，在保证抗量子安全性的同时降低连接建立时延。

\vspace{1em}

\noindent\textbf{摘要附图建议：}图1。

\end{document}
